% ============================================
% Visiting Assistant Professor Cover Letter – Middlebury College
% ============================================

\documentclass[11pt,letterpaper]{letter}

\usepackage{geometry}
\usepackage{microtype}
\usepackage[hidelinks]{hyperref}

\geometry{left=25mm,right=25mm,top=25mm,bottom=25mm}

\signature{Decory Edwards}
\address{
Department of Economics \\
Johns Hopkins University \\
Baltimore, MD 21218 \\
\texttt{dedwar65@jh.edu}
}

\begin{document}

\begin{letter}{Search Committee \\
Department of Economics \\
Middlebury College \\
Middlebury, VT \\ USA}

\opening{Dear Members of the Search Committee,}

I am writing to apply for the Visiting Assistant Professor position in the Department of Economics at Middlebury College for the 2026--2028 academic years. I am a Ph.D.\ candidate in Economics at Johns Hopkins University, advised by Professors Christopher D.~Carroll, Nicholas W.~Papageorge, and Stelios Fourakis, and I expect to receive my degree in May 2026. My research fields are macroeconomics, wealth and income dynamics, and computational economics.

My research agenda focuses on understanding the determinants of wealth inequality and the mechanisms through which heterogeneity in returns to wealth shapes aggregate outcomes. In my job-market paper, I estimate the degree of heterogeneity in individual portfolio returns using micro-data from the Survey of Consumer Finances and simulate a structural model in which households face idiosyncratic return risk. I show that allowing for persistent heterogeneity in returns substantially improves the model’s ability to match the empirical distribution of wealth, offering a tractable way to connect micro-level return dispersion to macroeconomic inequality.

In a second project, I use the Health and Retirement Study to examine the relationship between interpersonal trust and portfolio performance. Building on the literature that finds a hump-shaped relationship between trust and income, I show that a similar relationship holds for individual returns to wealth measured in the HRS data, highlighting a behavioral channel through which social attitudes shape saving and investment outcomes. This work also provides new estimates of the persistent component of returns to net worth, which motivate the calibration of heterogeneity in my structural model.

Middlebury College is an especially strong fit for me as a visiting faculty member. The department’s emphasis on excellence in undergraduate teaching, careful development of economic intuition, and close faculty–student interaction aligns closely with my own approach to economics instruction. I am particularly drawn to Middlebury’s liberal arts environment, where rigorous economic reasoning is combined with broad intellectual engagement and where visiting faculty play an integral role in the core curriculum.

\textbf{Teaching.} I have experience teaching and supporting courses in \textit{Principles of Macroeconomics}, \textit{Intermediate Macroeconomic Theory}, and applied macroeconomic electives. My teaching philosophy emphasizes clarity, structure, and economic intuition, with frequent use of data and real-world examples to motivate formal models. In an undergraduate-focused setting like Middlebury, I prioritize active learning, clear scaffolding of quantitative concepts, and regular feedback to support students with diverse academic backgrounds. I would be well prepared to teach introductory and intermediate microeconomics or macroeconomics, as well as statistics or regression courses, and I would be enthusiastic about contributing to the department’s core teaching mission during the visiting appointment.

I would welcome the opportunity to contribute to Middlebury’s academic community during this two-year visiting position and to support the department’s teaching needs while remaining an active and engaged scholar.

Thank you very much for your time and consideration. I would be delighted to discuss my application further.

\closing{Sincerely,}

\end{letter}
\end{document}
